% Part I, chapter 3: Where files are read and written.
\chapter{Where files are read and written}
\label{chap:files}

\wfdbx{} reads and writes a file exactly where its name says, and nowhere
else. A full path, a \code{\textasciitilde} path or a URL is used as given;
any other name is read from the \emph{input directory} and written to the
\emph{output directory}, which are both the current directory unless they
are set. Nothing is searched, and URLs are opened only when network access
has been turned on. This differs from the C library, which searched a list
of places, the WFDB path, for every file it read (see
\cref{sec:named-files} for the reasons); legacy mode restores the search
(\cref{sec:legacy}).

\section{The input and output directories}
\label{sec:directories}

\Cref{tab:directories} shows the three settings that decide where files
are read and written. Each can be set by a library function, which has the
same name in every language, and by an environment variable; a call of the
function overrides the environment variable.

\begin{table}[htbp]
\caption{The settings that locate files.}
\label{tab:directories}
\begin{tabularx}{\linewidth}{@{}P{3.1cm}LP{2.6cm}P{2.2cm}@{}}
\toprule
{\tabhead Setting} & {\tabhead Functions} & {\tabhead Environment variable}
  & {\tabhead Default} \\
\midrule
input directory (a directory, a \code{\textasciitilde} path, or a URL) &
  \code{wfdbx\_set\_input\_dir}, \code{wfdbx\_get\_input\_dir} &
  \code{WFDBX\_INPUT\_DIR} & the current directory \\
output directory (a local directory or a \code{\textasciitilde} path) &
  \code{wfdbx\_set\_output\_dir}, \code{wfdbx\_get\_output\_dir} &
  \code{WFDBX\_OUTPUT\_DIR} & the current directory \\
network access & \code{wfdbx\_set\_network}, \code{wfdbx\_get\_network} &
  \code{WFDBX\_NETWORK} & off \\
\bottomrule
\end{tabularx}
\end{table}

\begin{itemize}
\item \code{wfdbx\_set\_input\_dir(d)} sets the input directory, and returns
  the previous one (or none, if none was set). \code{d} is a local
  directory, a \code{\textasciitilde} path, or the URL of a directory, such
  as \code{https://physionet.org/files/mitdb/1.0.0/}; reading from a URL also
  needs network access. None (null, \code{NULL}, \ldots, depending on the
  language) or an empty string restores the default: the value of
  \code{WFDBX\_INPUT\_DIR} if it is set, and otherwise the current directory.
  \code{wfdbx\_get\_input\_dir()} returns the input directory, or none.
\item \code{wfdbx\_set\_output\_dir(d)} sets the output directory in the
  same way, and returns the previous one; the default is the value of
  \code{WFDBX\_OUTPUT\_DIR}, or the current directory. The output directory
  must be local: a name that begins with a URL scheme, such as
  \code{https://} or \code{ftp://} (given to the function, or in the
  environment variable), is refused with the message
  \code{wfdbx\_set\_output\_dir: the output directory must be a local
  directory}, and the setting does not change.
\item \code{wfdbx\_set\_network(on)} turns network access on (1, or true) or
  off (0, or false), and returns the previous setting (1 or 0); none or a
  negative number restores the default: on if \code{WFDBX\_NETWORK} is set
  to a nonzero number (such as \code{1}), off otherwise.
  \code{wfdbx\_get\_network()} returns 1 or 0.
\end{itemize}
In Kotlin, these functions are also available under camelCase names
(\code{wfdbxSetInputDir}, \ldots). The settings apply to the whole library:
to the high-level functions, the library functions and the tools, which use
the environment variables, also for the files that they open themselves
(\cref{sec:tool-files}). \code{wfdbquit} does not reset them.

The high-level functions also take these settings as options, which apply
for the duration of one call, after which the previous settings are
restored: the readers take an input directory (\code{input\_dir}) and a
network permission (\code{network}), and the writers an output directory
(\code{dir}). Each of these options is optional; \code{input\_dir} applies
to relative record names only, and \code{dir} to plain record names
(\cref{sec:record-paths}). In Kotlin the options are \code{inputDir},
\code{network} and \code{dir}.

\paragraph{The C functions setwfdb and getwfdb.} In \wfdbx{},
\code{setwfdb} and \code{getwfdb} set and return the input directory.
\code{setwfdb} splits its argument as the C library splits a WFDB path (at
spaces, semicolons, and colons that are not part of a URL), and accepts a
single directory (\code{setwfdb(".")}); a path of several directories is an
error (\code{setwfdb: WFDB-x reads from a single input directory; use
wfdbx\_set\_input\_dir (or legacy mode for the WFDB search path)}), and
nothing changes. A directory whose name contains a space must therefore be
set with \code{wfdbx\_set\_input\_dir}, which takes the name as it is.
\code{getwfdb} returns \code{.} when no input directory is set.
\code{resetwfdb} restores the directory set by the last \code{setwfdb}.

\section{Record names}
\label{sec:record-paths}

A record name is given without the extension \file{.hea} of its header
(\code{100s}, not \code{100s.hea}; a trailing \file{.hea} is removed by the
functions that read a record and by \code{write\_record}). It is one of:
\begin{itemize}
\item a full path, absolute (\code{/data/mitdb/100}) or in the home
  directory (\code{\textasciitilde/Documents/100s}), used as given;
\item a relative path (\code{100s}, \code{data/100s}), read from the input
  directory if one is set, otherwise from the current directory, and
  written to the output directory if one is set, otherwise to the current
  directory (missing directories are created);
\item a URL (\code{https://physionet.org/files/mitdb/1.0.0/100}), read only
  if network access is on (\cref{sec:network}); a URL cannot be written.
  Only \code{http://} and \code{https://} URLs can be read: a name with
  another URL scheme, such as \code{ftp://} or \code{sftp://}, is an error,
  and is never read as a local file name.
\end{itemize}
A writer given a record name with a path (\code{out/rec}) creates the files
in that directory, and the header names the record and its signal file
without the path (\code{rec}, \code{rec.dat}), so that the record opens
under any path that leads to it (\cref{sec:resolve}). With the
\code{dir} option of a writer, the record name must be a plain name
(\code{rec}), without a directory: the files are created in \code{dir},
which must exist, and the record opens from that directory under its plain
name. The characters allowed in record names are given in
\cref{sec:record-names}.

\section{How a name is resolved}
\label{sec:resolve}

The name of a file is the record name followed by a dot and the file type
(\file{100s.hea}, \file{100s.atr}), or a file name given by itself (a signal
file named in a header, a calibration file). It designates exactly one
location, as \cref{tab:resolve} shows.

\begin{table}[htbp]
\caption{Where a name is read from and written to.}
\label{tab:resolve}
\begin{tabularx}{\linewidth}{@{}P{4.4cm}LL@{}}
\toprule
{\tabhead The name is} & {\tabhead Read from} & {\tabhead Written to} \\
\midrule
an absolute path: \code{/data/100}; on MS-Windows also
  \code{C:\textbackslash{}data\textbackslash{}100}, \code{C:/data/100} or
  \code{\textbackslash\textbackslash{}server\textbackslash{}share\textbackslash{}100} &
  that path & that path \\
a home path: \code{\textasciitilde/data/100}, \code{\textasciitilde
  user/data/100} & the home directory (expanded by \wfdbx{}) & the same \\
a URL: \code{http://\ldots} or \code{https://\ldots} & that URL, if network
  access is on & (not possible) \\
a name with another URL scheme: \code{ftp://\ldots}, \code{sftp://\ldots},
  \code{file://\ldots} & (not possible: an error) & (not possible) \\
\code{-} (as a record name) & the standard input & the standard output \\
anything else: \code{100}, \code{mitdb/100}, \code{../other/100} & the input
  directory followed by the name & the output directory followed by the
  name \\
\bottomrule
\end{tabularx}
\end{table}

\begin{itemize}
\item A leading \code{\textasciitilde} or \code{\textasciitilde/} is
  replaced by the home directory (the environment variable \code{HOME}, or
  the home directory of the current user if it is not set; on MS-Windows,
  \code{USERPROFILE}), and, on POSIX systems, \code{\textasciitilde
  user} by the home directory of that user. The input and output
  directories may themselves be \code{\textasciitilde} paths.
\item The input or output directory and the name are joined with one
  directory separator (\code{/} in a URL). \code{..} means one level up;
  the C library refused every name that contains \code{..}.
\item Nothing else is tried: a file that is not at that location is not
  found, and the operation fails (for example, \code{init: can't open header
  for record 100}). The only other name tried at the same location is the
  one whose file type is cut to its first three characters, as in C.
\item When an output file cannot be created because a directory of its path
  does not exist, the missing directories are created, as in C.
\end{itemize}
A record written under a path, such as \code{out/rec}, gets its files in
\file{out/} (under the output directory, if it is relative). Only the
record's own name, the last component, must be made of the characters of
record names (\cref{sec:record-names}); the directories may contain any
character. The header names the record, its signal files and its segments
relative to its own directory (\code{rec}, \code{rec.dat}), which is where
they are read from; the C library wrote the names as given, with their path
(\code{out/rec} and \code{out/rec.dat} in \file{out/rec.hea}), as legacy
mode does, so that such a record could not be read back by finding its
files relative to its header.

\subsection{Files named in a header}
\label{sec:header-files}

The files that a header names, the signal files of a record and the headers
of the segments of a multi-segment record, are found relative to the
directory of the header, as they are written in it (\code{100s.dat}, or
\code{data/100s.dat}): on disk, or at the URL from which the header was read.
A file missing there is an error, not a search elsewhere. Such a name must
stay inside the header's directory: a header that names a file with a
\code{..} component, an absolute path, a \code{\textasciitilde} path or a
URL (or, on MS-Windows, a drive letter) is rejected, with the message
\begin{lstlisting}[style=text]
init: file NAME named in the header of record REC is not in the header's directory
\end{lstlisting}
The signal file name \code{-} (the standard input) is accepted, and the null
segment \code{\textasciitilde} of a multi-segment record keeps its meaning.
For a header read from the standard input, which has no directory, the
files it names are resolved as any relative name, under the input
directory. This rule keeps a downloaded or crafted header from making the
library read other local files or contact another server.

\subsection{Files that a tool opens itself}
\label{sec:tool-files}

A few tools open a file themselves, by the name given on the command line,
instead of through the library. Such a file is found, or created, exactly
where the library finds, or creates, a file of that name: a relative name
under the input directory (a file that the tool reads) or the output
directory (a file that it creates, whose missing directories under the
output directory are created), a \code{\textasciitilde} name expanded, and
an absolute name as given; \code{-}, where the tool takes it for the
standard input or output, keeps that meaning. A URL is not accepted for
these files, also when the input directory is a URL: the tool reports that
it cannot read, or create, the file, as for a file that does not exist.
Messages name the file as it was given. \Cref{tab:tool-files} lists these
files.

\begin{table}[htbp]
\caption{The files that the tools open themselves.}
\label{tab:tool-files}
\begin{tabularx}{\linewidth}{@{}P{3.4cm}LL@{}}
\toprule
{\tabhead Tool} & {\tabhead Reads} & {\tabhead Creates} \\
\midrule
\cmd{wrsamp}, \cmd{edf2mit}, \cmd{wav2mit}, \cmd{a2m}, \cmd{ad2m} & the
  \opt{-i} file & \\
\cmd{ahaecg2mit} & the input files (\file{.ecg}, \file{.txt}) & \\
\cmd{xform} & the script of \opt{-S} & \\
\cmd{wfdbdesc} & the first signal file, to compute the length of a record
  whose header gives none: in the header's directory, where the library
  reads it & \\
\cmd{mit2wav} & & the WAV file (\opt{-o}), completed where the library
  wrote its samples \\
\cmd{mit2edf} & & the EDF file (\opt{-o}, or \textit{RECORD}\file{.edf}) \\
\cmd{snip} & & the master header of a multi-segment excerpt (\opt{-m}) \\
\cmd{parsescp} & & the \file{.des}, \file{.ecg}, \file{.key} and
  \file{.txt} files \\
\bottomrule
\end{tabularx}
\end{table}

The library functions \code{wfdbx\_input\_path} and
\code{wfdbx\_output\_path} (\cref{sec:library-additions}) return these
locations, for any program that opens a file itself. The C tools opened
these files relative to the current directory, so that, with an input or
output directory set, a tool read its records from one place and its own
files from another (\cmd{mit2wav} left the WAV header in the current
directory and the samples, without a header, in the output directory).
Without an input or output directory, nothing changes, except that a
\code{\textasciitilde} name is expanded and that \cmd{wfdbdesc} looks for
the signal file of a record in another directory in that directory.

\subsection{Legacy mode}

In legacy mode (\cref{sec:legacy}), files are found as by the original
reference build of the C library, which has no network support: an input
file is searched for in each entry of the WFDB path in turn, given by the
\code{WFDB} environment variable or \code{setwfdb} (by default
\code{. /usr/local/database http://physionet.org/physiobank/database}), with
the \code{\%r} substitution of the record name and \code{@file} indirect
paths; absolute names are joined to the entries of the path (an empty entry,
as in \code{WFDB=";"}, makes names be used as given); names that contain
\code{..} are refused, and URLs are not opened. The input directory and
\code{WFDBX\_NETWORK} are not used. The output directory is used, as a plain
prefix of every output name, absolute names included. The tools open the
files that they read themselves (\cref{sec:tool-files}) as named, relative
to the current directory, as the original tools do, and the output
directory, if one is set, is a plain prefix of the names of the files that
they create themselves, as of every other output file.

\section{Network access}
\label{sec:network}

Network access is off by default. A URL is opened only when it has been
turned on, by \code{wfdbx\_set\_network(1)}, by \code{WFDBX\_NETWORK=1}, or
by the \code{network} option of a high-level reader. Otherwise the library
says once, the first time it is asked to open a URL:
\begin{lstlisting}[style=text]
WFDB-x: network access is disabled; set WFDBX_NETWORK=1 or call wfdbx_set_network(1) to open URL
\end{lstlisting}
and the file is not found. The files of PhysioNet are reached in the same
way as any other remote files: by setting the input directory to the URL of
the database, \code{https://\allowbreak physionet.org/\allowbreak files/}\emph{database}\code{/}%
\emph{version}\code{/}, and turning network access on. No URL is built in:
\wfdbx{} never contacts a server it was not told to contact. For example, to
print the first 10 seconds of record 100 of the MIT-BIH Arrhythmia Database:
\begin{lstlisting}[style=shell]
export WFDBX_INPUT_DIR=https://physionet.org/files/mitdb/1.0.0/
export WFDBX_NETWORK=1
rdsamp -r 100 -t 10
\end{lstlisting}
or, in Python:
\begin{lstlisting}[language=Python]
import wfdbx
wfdbx.wfdbx_set_network(1)
wfdbx.wfdbx_set_input_dir("https://physionet.org/files/mitdb/1.0.0/")
rec = wfdbx.read_record("100", sampto="0:10")
\end{lstlisting}
or, for one call only:
\begin{lstlisting}[language=Python]
rec = wfdbx.read_record("100", input_dir="https://physionet.org/files/mitdb/1.0.0/",
                        network=True)
\end{lstlisting}

Remote files can only be read. Each is downloaded in full the first time it
is opened, and kept for the rest of the process; a download that fails is
not kept, so the next time the file is opened it is requested again (a
server that was unavailable may have recovered). Redirects are followed, and
an HTTP status of 400 or more is a failure. A download gives up if
connecting takes 30 seconds, or if no data arrive for 30 seconds; there is
no limit on its total time. Each language downloads with its own means:
\begin{itemize}
\item Python: its standard \code{urllib};
\item C++ and Rust: the system \cmd{curl} program, which they run without a
  shell (a missing \cmd{curl} is reported when a URL is opened);
\item Julia: its standard \code{Downloads};
\item R and MATLAB/Octave: \cmd{curl} when it is installed, and otherwise
  R's \code{url()} or MATLAB's \code{webread} and \code{websave};
\item Kotlin: \code{java.net.HttpURLConnection} on the JVM and Android, and
  \code{NSURLSession} on iOS.
\end{itemize}

Only \code{http://} and \code{https://} URLs can be read. A name that begins
with another URL scheme (\code{ftp://}, \code{ftps://}, \code{sftp://},
\code{file://}, \ldots) is refused, with or without network access, with a
message the first time:
\begin{lstlisting}[style=shell]
WFDB-x: can't open ftp://example.org/db/100.hea: only http:// and https:// URLs
can be read (mount the server as a local folder, or download the files)
\end{lstlisting}
(on one line). The same holds for an input directory with such a scheme,
a file named in a header, and the name of an output file; the files that a
tool reads or creates itself are refused as URLs are
(\cref{sec:tool-files}), and an output directory with a URL scheme is
refused (\cref{sec:directories}). Such a name is never read as a local file
name (the directory \file{ftp:} is never looked in). A scheme is a letter
followed by at least one letter, digit, \code{+}, \code{-} or \code{.},
then \code{://}, so a one-letter prefix, as in \code{c://x} (like a drive
letter), is not one. The records of a server reached by FTP or SFTP can be
read once the server is mounted as a local folder, by their local paths, or
after the files are downloaded:
\begin{lstlisting}[style=shell]
sshfs user@example.org:/data/db ~/mnt/db       # SFTP (or: rclone mount)
rdsamp -r ~/mnt/db/100 -t 10
\end{lstlisting}
In legacy mode, every name containing \code{://} is refused, without a
message, as by the original library built without network support.

A header read from a URL can be read but not changed by the metadata
functions, except into a local directory (\cref{sec:metadata-file}).

\section{The standard input and output}
\label{sec:stdio}

The record name \code{-} stands for the standard input, when a file is read,
and for the standard output, when one is written, whatever the type of the
file: a header, an annotation file or a signal file. A signal file named
\code{-} in a header is the standard input (or output). The tools follow the
C tools in this, so that, for example, \cmd{wrsamp} reads its text input from
the standard input unless \opt{-i} is given. The metadata functions refuse
the record name \code{-}.

\section{Environment variables}
\label{sec:environment}

\Cref{tab:environment} lists every environment variable that \wfdbx{} reads.
The library reads each one when it needs it.

\begin{xltabular}{\linewidth}{@{}P{3.9cm}L@{}}
\caption{The environment variables used by \wfdbx{}.}
\label{tab:environment}\\
\toprule
{\tabhead Variable} & {\tabhead Meaning} \\
\midrule
\endfirsthead
\toprule
{\tabhead Variable} & {\tabhead Meaning} \\
\midrule
\endhead
\bottomrule
\endfoot
\code{WFDBX\_INPUT\_DIR} & The input directory (\cref{sec:directories}):
  a local directory, a \code{\textasciitilde} path or a URL. Not used in
  legacy mode. \\
\code{WFDBX\_OUTPUT\_DIR} & The output directory: a local directory or a
  \code{\textasciitilde} path. \\
\code{WFDBX\_NETWORK} & Network access is on if this is a nonzero number
  (\code{WFDBX\_NETWORK=1}). Not used in legacy mode. \\
\code{WFDB} & The WFDB path, used only in legacy mode. Outside legacy mode
  it is ignored, and if it is set the library says so once:
  \code{WFDB-x: the WFDB environment variable is ignored; set
  WFDBX\_INPUT\_DIR instead (or use legacy mode)}. In legacy mode,
  \code{setwfdb} sets it (and sets \code{WFDBCAL}, \code{WFDBANNSORT} and
  \code{WFDBGVMODE} to their defaults if they are not set), as the C library
  does, so that the programs it starts inherit them. \\
\code{WFDBANNSORT} & If it is \code{0}, an output annotation file that is
  not in canonical order is not sorted when it is closed; the library then
  prints the advice \code{Use the command: sortann \ldots} (\cref{sec:sorting}).
  The \code{sort} option of the high-level \code{write\_annotations} decides
  for its call, whatever this variable says. (\code{WFDBNOSORT}, named by
  an old manual page of \cmd{sortann}, is not used, as in C.) \\
\code{WFDBCAL} & The name of the calibration file that \code{calopen} reads
  when it is given no name; \file{wfdbcal} if it is not set
  (\cref{sec:calibration}). \\
\code{WFDBGVMODE} & The \code{getvec} mode used when \code{setgvmode} is
  given a negative mode, and by \cmd{sigamp} when its options do not set the
  mode: \code{0} for low resolution, \code{1} for high resolution, plus
  \code{2} to replace invalid samples by the previous valid ones
  (\cref{sec:frames}). \\
\code{WFDB\_FLAC\_CHECK\_MD5} & If it is set to a nonempty value, the MD5
  signature of each FLAC signal file that is read is verified when the file
  is closed; a mismatch gives \code{isigclose: warning: incorrect MD5 hash
  in FLAC input}. \\
\code{HOME}, \code{USERPROFILE} & The home directory, for names that begin
  with \code{\textasciitilde} (\code{USERPROFILE} on MS-Windows). \\
\code{PATH} & Where the system \cmd{curl} program is found (C++, Rust, and
  R and MATLAB/Octave), and, in legacy mode, the \cmd{sortann} program that
  the library runs to sort an annotation file (C++ and Rust look first in
  the directory of the running program). \\
\code{TZ} & On POSIX systems, the time zone of the local time, which the
  library uses where the C library does: for the current date and time
  (\code{setbasetime} with an empty string, and \cmd{mit2edf} for a record
  without a start time). \\
\code{WFDBX\_DEBUG} & (R and Kotlin tools.) If it is set, a tool that stops
  on an internal error (reported as \code{wfdbx: error: \ldots}) also prints
  the call stack. \\
\code{WFDBX\_STDIN} & (R and MATLAB/Octave.) The name of a file that is
  read as the standard input: in R, for the tests; in MATLAB/Octave, also
  set by the MATLAB driver \file{matlab/bin/wfdbx\_matlab.sh}, since MATLAB
  cannot read its standard input. \\
\code{WFDBX\_STDOUT}, \code{WFDBX\_STDERR} & (MATLAB/Octave.) Files to which
  the standard output and error are written, set by the MATLAB driver,
  which copies them to the real standard output and error when MATLAB
  exits. \\
\end{xltabular}

No environment variable turns legacy mode on: it is turned on for a tool by
its \opt{--legacy} option, and in a program by \code{wfdbx\_set\_legacy(1)}
or the \code{legacy} option of a high-level function (\cref{sec:legacy}).
The other variables that the C library reads are not used by \wfdbx{}: those
of its libcurl code (\code{WFDB\_PAGESIZE}, \code{WFDBPASSWORD},
\code{CURL\_CA\_BUNDLE}, \code{WFDB\_NET\_DEBUG}) and those that tune
libFLAC's encoder (\code{WFDB\_FLAC\_COMPRESSION\_LEVEL},
\code{WFDB\_FLAC\_BLOCK\_SIZE}, and the other \code{WFDB\_FLAC\_\ldots}
variables).

The conformance suite runs every case with \code{LC\_ALL=C} and \code{TZ=UTC},
and without the variables \code{WFDB}, \code{WFDBCAL}, \code{WFDBANNSORT},
\code{WFDBGVMODE}, \code{WFDBX\_INPUT\_DIR}, \code{WFDBX\_OUTPUT\_DIR} and
\code{WFDBX\_NETWORK}, so that its results do not depend on them.
